%==================== CHAPTER 2 ====================
\ChapterPage{Study Of the System}

\section{Literature Survey}

Building an agentic AI platform for business functions means drawing on work from several overlapping areas of research. We looked at how multi-agent frameworks are structured, how AI has been applied to hiring and whether it can be trusted to do so fairly, how marketing content generation and hallucination reduction have been approached, and what the state of the art says about tool use, privacy and governance. The picture that emerges is one of genuine progress alongside well-documented gaps, and those gaps are what this project tries to close.

\vspace{10pt}
\noindent\textbf{Multi-Agent Orchestration.}
In AutoGen, Wu et al. \cite{autogen} proposed a framework that enables users to write applications using multiple agents, with each agent using an LLM, a human, or an external tool. Guo et al. \cite{guo} followed with a very general overview of multi-agent systems based on LLM, including the profile of agents, agent communication, and the development of agent capabilities over time. Yao et al.\cite{react} found that the alternating strategy of reasoning steps and tool actions yields better results than either. None of these can provide an out-of-the-box approach to a domain-specific business workflow. They provide us with building blocks, but it's going to take some work to design the plumbing for rubric-based scoring, as well as human approval gates and decision-level audit logs.
\vspace{10pt}
\noindent\textbf{AI-Powered Recruitment and Bias.}
The recruitment literature makes a sobering case for caution. Gan et al.\ \cite{gan} showed that an LLM-based pipeline can screen resumes eleven times faster than a human and reach a respectable F1 score, but their system was tested only on IT-sector data and could not generate rubrics from a job description or conduct an interview. Wilson and Caliskan \cite{wilson} found that embedding-based screening favoured White-associated names in 85 out of every 100 cases and that Black male candidates were disadvantaged in up to 100\% of test cases. Mu et al.\ \cite{mu} then showed that adversarial content hidden inside a resume could manipulate LLM screening more than 80\% of the time, and that the best available defences reduced this by only around a quarter. Stuart et al.\ \cite{stuart} took things further by building an AI interviewer that updates its belief about a candidate's competencies after each answer, recovering accurate ability profiles with 76.1\% accuracy. Together these papers describe a picture where AI recruitment tools can save time but cannot yet be trusted without human oversight, bias monitoring and adversarial scanning at every step.

\vspace{10pt}
\noindent\textbf{Content Generation and Marketing.}
Liu et al.\ \cite{marketingfm} demonstrated that retrieval-augmented generation can produce keyword-specific ad copy that outperforms template-based content in online A/B tests, with an automated evaluation approach that agrees with human reviewers nearly 90\% of the time. The limitation is scope: the system was designed for search advertising rather than for brand-consistent posts across multiple social channels, and it does not include an approval step before content goes live. Our marketing workflow extends the RAG-based generation idea to channel-specific branded content and adds an explicit queue where a human must approve every draft before it is published.

\vspace{10pt}
\noindent\textbf{Retrieval-Augmented Generation and Hallucination Reduction.}
To evaluate whether to retrieve or not, and to judge whether the retrieved facts are accurate or not, Asai et al.\ \cite{selfrag} trained a model to do both using special reflection tokens, achieving better factuality scores than baseline retrieval-augmented approaches. In another approach, Dhuliawala et al. \cite{cove} drew the model to draft an answer, created verification questions and then refined the answer according to the verification results, resulting in fewer hallucinations on some benchmarks. Both of these ideas are actually flawed in practice; either they require custom training for the machine to learn the new information or they only work if the information is something the machine already knows. Our podcast research agent takes the retrieve-then-verify concept, and implements it with prompting: Any statements that are made are made from a public source, and marked as a "verified fact", an "allegation" or an "opinion" before being directed to the host.
\vspace{10pt}
\noindent\textbf{Tool Use and Enterprise API Integration.}
Schick et al.\ \cite{toolformer} showed that a language model can learn which API to call and when, purely from a small number of demonstrations, without supervised training. Patil et al.\ \cite{gorilla} went further, fine-tuning a model to write accurate API calls even when the documentation changes after training. Both results are impressive, but neither system can chain multiple authenticated calls, handle rate limits or decide whether a particular write operation should require human approval before it goes through. Enterprise integrations with Zoho, Workday and Greenhouse require all of these things, which is why our connector layer adds OAuth scopes, retries, idempotency keys and a mandatory approval step for any call that writes data.

\vspace{10pt}
\noindent\textbf{Privacy and Governance.}
Lukas et al.\ \cite{lukas} conducted an empirical study of the extent to which a language model can expose PII sequences and discovered that even with techniques such as differential privacy, approximately 3\% of the PII sequence is still exposed. M\"{o}kander et al. \cite{mokander} came up with a governance blueprint for LLMs at organisational, model and application levels but did not include how to do application-level auditing in software in the paper. We take a proactive approach to PII risk by masking all PII data in a candidate's response before it is sent to an LLM and limiting the data that is logged. In this case, application audit layer is not a design concept but a real implementation: every prompt, tool call, source, draft version, reviewer decision and final output are recorded and tracked.

\clearpage
{\footnotesize
\renewcommand{\arraystretch}{1.35}
\begin{longtable}{|>{\raggedright\arraybackslash}p{0.75cm}|>{\raggedright\arraybackslash}p{2.6cm}|>{\raggedright\arraybackslash}p{3.75cm}|>{\raggedright\arraybackslash}p{3.0cm}|>{\raggedright\arraybackslash}p{3.4cm}|}
\caption{Summary of Related Research Papers}\label{tab:lit}\\
\hline
\textbf{S.no} & \textbf{Paper} & \textbf{Main Idea / Objective} & \textbf{Tools / Tech Used} & \textbf{Gaps Identified}\\ \hline
\endfirsthead
\hline
\textbf{S.no} & \textbf{Paper} & \textbf{Main Idea / Objective} & \textbf{Tools / Tech Used} & \textbf{Gaps Identified}\\ \hline
\endhead
\endlastfoot
1 & AutoGen (2024) & Framework for building applications from multiple conversable agents that combine LLMs, humans and tools. & Multi-agent conversation framework; human-in-the-loop support. & General-purpose; no HR or marketing workflow, rubric scoring, approval gates or audit design.\\ \hline
2 & LLM Based Multi-Agents: A Survey (2024) & Reviews domains, agent profiling, communication, capability growth and benchmarks of LLM multi-agent systems. & Systematic literature review; open-source tracking repository. & Survey only; no concrete business-function system, governance or integration.\\ \hline
3 & ReAct (2023) & Interleaves reasoning traces with tool actions to reduce hallucination and improve decision-making. & Few-shot prompting; Wikipedia API; HotpotQA, FEVER, ALFWorld, WebShop. & Single agent with read-only tools; no approval or audit for actions with real effects.\\ \hline
4 & LLM Agents in Recruitment (2024) & Agent pipeline that classifies, grades, summarises and selects resumes. & Fine-tuned LLaMA2; GPT-3.5 and GPT-4 baselines; BERT, T5. & IT resumes only; labels mostly from GPT-4; no JD-to-rubric generation or interview stage.\\ \hline
5 & Bias in Resume Screening via LM Retrieval (2024) & Audits embedding-based screening for gender, race and intersectional bias. & Massive Text Embedding models; retrieval-based audit on 500+ resumes. & Measures bias; proposes no bias-aware screening or human review step.\\ \hline
6 & Adversarial Vulnerabilities in Resume Screening (2025) & Shows hidden instructions and invisible content can manipulate LLM screening. & 12 LLM variants; prompt defence; FIDS (LoRA). & Defences cut attacks by only 10--26\%; no pipeline-level scanning or human checkpoint.\\ \hline
7 & Beyond the Resum\'e (2026) & Rubric-aware LLM interviewer that updates beliefs about competencies after each answer. & GPT-5; belief updating; metamorphic testing; Chainlit. & Simulated interviews only; three-level rubrics; no rubric generation, ATS link or HR approval.\\ \hline
8 & LLMs for Marketing Content Generation (2025) & Generates keyword-specific ad copy and evaluates it at scale. & RAG; LLM-as-a-Judge; rule-based metrics; A/B testing. & Ad copy only; no multi-channel brand posts or approval workflow.\\ \hline
9 & Self-RAG (2024) & Model learns when to retrieve and critiques its own output with reflection tokens. & Fine-tuned Llama2 7B/13B; retriever; reflection tokens. & Needs custom training; not evaluated on high-risk research about real people.\\ \hline
10 & Chain-of-Verification (2024) & Drafts, plans and answers verification questions, then revises to reduce hallucination. & Multi-stage prompting; factored verification. & Uses only the model's own knowledge; no external sources or claim labels.\\ \hline
11 & Toolformer (2023) & LLM teaches itself when and how to call external APIs. & GPT-J 6.7B; self-supervised API-call annotation. & One call per example; no tool chaining, authentication or write actions.\\ \hline
12 & Gorilla (2024) & Fine-tuned LLM writes accurate API calls with retrieval over documentation. & LLaMA-7B; BM25 and GPT-Index retrievers; APIBench; AST matching. & Public ML APIs only; ignores permissions, PII and approval of write operations.\\ \hline
13 & PII Leakage in Language Models (2023) & Defines and measures PII extraction, reconstruction and inference attacks. & GPT-2 fine-tuning; black-box attacks; differential privacy. & Covers training-data leakage; not PII in run-time prompts, logs and tools.\\ \hline
14 & Auditing LLMs: A Three-Layered Approach (2024) & Blueprint of governance, model and application audits for LLMs. & Conceptual governance framework. & Conceptual; no implementation of application-level audit trails.\\ \hline
\end{longtable}
}

\clearpage
\section{Existing Systems}

There are a number of commercial tools currently addressing the overlapping portions of the problem, but nothing available that could integrate HR screening with marketing content creation, podcast research verification and enterprise automation in one managed workflow. While all have their strengths in recruitment (HireVue, Greenhouse, and Zoho Recruit), marketing (Jasper), general research (Perplexity Deep Research), and automation (n8n), they also have weaknesses in other critical functions including rubric generation, claims supported by evidence, source accountability, auditability, approval gates, and end-to-end business workflow orchestration. To put it simply, the existing market is filled with point solutions for domains, but not a single agentic platform for business operations that is policy-aware.
\clearpage
